\chapter{Gramiany obserwowalności i realizacje minimalne}
\label{ch:lin27-realizations}
\index{gramian!obserwowalności}
\index{realizacja!minimalna}
\index{macierz Hankela}
\index{układ liniowy!realizacja}
\index{realizacja!zrównoważona}

\noindent\textbf{Cel rozdziału.}
W Rozdziale~\ref{ch:lin25-control}
zdefiniowaliśmy przestrzeń osiągalną
i nieobserwowalną, wprowadziliśmy
kryteria Kalmana, gramian sterowalności
i rozkład na odpowiednie podprzestrzenie.
Obecnie zbadamy dokładniej ilościowy
sens \emph{gramianu obserwowalności}
i skonstruujemy układ o najmniejszej
liczbie zmiennych stanu, który zachowuje
wszystkie odpowiedzi wejście--wyjście.
Podamy algebraiczne kryterium
minimalności przez macierz Hankela
oraz wyjaśnimy, w jakim sensie
minimalna realizacja jest jednoznaczna.
W stabilnych układach skonstruujemy
również współrzędne, w których
oba gramiany są jednocześnie diagonalne
i równe.

Niech \(n,m,p\ge1\) i niech
\(A\in\operatorname{Mat}_{n\times n}(\mathbb C)\),
\(B\in\operatorname{Mat}_{n\times m}(\mathbb C)\),
\(C\in\operatorname{Mat}_{p\times n}(\mathbb C)\).
Przestrzeń stanów oznaczamy
przez \(V=\mathbb C^n\).
Rozważamy układ dyskretny
\(x_{k+1}=Ax_k+Bu_k\), \(y_k=Cx_k\),
oraz ciągły
\(x'=Ax+Bu\), \(y=Cx\).
W obu przypadkach
obserwacje i sterowania
należą do przestrzeni
odpowiednio \(\mathbb C^p\)
i \(\mathbb C^m\).
Przez \(T^*\) oznaczamy
sprzężenie hermitowskie macierzy.
Oznaczenia
\(\mathcal R(A,B)\)
i \(\mathcal N(C,A)\)
pochodzą z
Definicji~\ref{def:lin25-kalman}
i~\ref{def:lin25-observable}.

\section{Skończone gramiany obserwowalności}
\label{sec:lin27-finite}

\begin{definicja}[Dyskretny gramian obserwowalności]
\label{def:lin27-obs-discrete}
Dla każdego całkowitego \(N\ge1\)
definiujemy
\[
 W_{\mathrm o,N}^{\mathrm d}
 =\sum_{k=0}^{N-1}(A^*)^kC^*CA^k.
\]
Jest to macierz stopnia \(n\),
zwana \emph{dyskretnym gramianem
obserwowalności na horyzoncie \(N\)}.
Zapiszmy też macierz obserwacji
\[
 \mathcal O_N(C,A)
 =\begin{bmatrix}
     C\\CA\\\vdots\\CA^{N-1}
   \end{bmatrix}.
\]
\end{definicja}

\begin{twierdzenie}[Energia obserwacji i odtwarzanie stanu]
\label{thm:lin27-discrete-gramian}
Dla wszystkich \(N\ge1\)
i \(x\in\mathbb C^n\) zachodzi
\[
 \begin{aligned}
 W_{\mathrm o,N}^{\mathrm d}
   &=\mathcal O_N(C,A)^*
      \mathcal O_N(C,A),\\
 x^*W_{\mathrm o,N}^{\mathrm d}x
   &=\sum_{k=0}^{N-1}\|CA^kx\|_2^2.
 \end{aligned}
\]
Macierz
\(W_{\mathrm o,N}^{\mathrm d}\)
jest hermitowska
i nieujemnie określona.
Dla \(N\ge n\) mamy
\[
 \ker W_{\mathrm o,N}^{\mathrm d}
 =\mathcal N(C,A).
\]
W szczególności para \((C,A)\)
jest obserwowalna dokładnie wtedy,
gdy \(W_{\mathrm o,n}^{\mathrm d}>0\).

Jeżeli \(W_{\mathrm o,N}^{\mathrm d}>0\),
to stan początkowy \(x_0\)
w układzie bez wymuszenia
\(x_{k+1}=Ax_k\), \(y_k=Cx_k\)
jest odtwarzany z obserwacji
\(y_0,\ldots,y_{N-1}\) przez
\[
 x_0=
 (W_{\mathrm o,N}^{\mathrm d})^{-1}
 \mathcal O_N(C,A)^*
 \begin{bmatrix}
  y_0\\y_1\\\vdots\\y_{N-1}
 \end{bmatrix}.
\]
\end{twierdzenie}
\begin{proof}
Mnożenie pionowej macierzy
\(\mathcal O_N\) przez jej sprzężenie
hermitowskie z lewej strony daje
sumę iloczynów
\((CA^k)^*(CA^k)\),
co jest dokładnie
\(W_{\mathrm o,N}^{\mathrm d}\).
Po pomnożeniu przez \(x^*,x\)
otrzymujemy sumę kwadratów norm.
Dowodzi to nieujemnej określoności
i hermitowskości.

Dla dowolnej macierzy \(M\)
równość \(M^*Mx=0\)
jest równoważna \(Mx=0\),
ponieważ
\(x^*M^*Mx=\|Mx\|_2^2\).
Stąd
\(\ker W_{\mathrm o,N}^{\mathrm d}
 =\ker\mathcal O_N\).
Dla \(N=n\) jest to z definicji
\(\mathcal N(C,A)\).
Jeśli \(N>n\), to
twierdzenie Cayleya--Hamiltona
gwarantuje, że każdy iloczyn
\(CA^k\) dla \(k\ge n\)
jest kombinacją wcześniejszych
\(C,CA,\ldots,CA^{n-1}\).
Nowe równania nie zmieniają
zatem jądra macierzy obserwacji.
Dodatnia określoność
gramianu jest równoważna
jego zerowemu jądru,
czyli obserwowalności
zgodnie z
Definicją~\ref{def:lin25-observable}.

Dla układu bez wymuszenia
\(y_k=CA^kx_0\).
Zestawiając obserwacje
w jedną kolumnę \(Y\),
otrzymujemy
\(Y=\mathcal O_Nx_0\).
Mnożenie przez
\(\mathcal O_N^*\) daje
\(\mathcal O_N^*Y
 =W_{\mathrm o,N}^{\mathrm d}x_0\).
Odwracalność gramianu
umożliwia rozwiązanie
tego równania i prowadzi
do wzoru z tezy.
\end{proof}

\begin{uwaga}[Znane sterowanie nie utrudnia odtworzenia]
\label{rem:lin27-known-input}
Gdy sterowania \(u_0,\ldots,u_{N-2}\)
są znane, iteracja równania stanu
z
Twierdzenia~\ref{thm:lin25-kalman}
daje
\[
 y_k=CA^kx_0+
       \sum_{j=0}^{k-1}
             CA^{k-1-j}Bu_j
       \quad(0\le k<N).
\]
Dla \(k=0\) sumę przyjmujemy
za zerową. Odejmuje się więc
znaną sumę od każdego \(y_k\)
i stosuje dokładnie ten sam
wzór na \(x_0\).
Obserwowalność jest własnością
\((C,A)\), a nie wybranego
ciągu znanych sterowań.
\end{uwaga}

\begin{definicja}[Ciągły gramian obserwowalności]
\label{def:lin27-obs-cont}
Dla każdego \(T>0\)
definiujemy
\[
 W_{\mathrm o,T}^{\mathrm c}
 =\int_0^T
    e^{tA^*}C^*Ce^{tA}
   \,\mathrm dt.
\]
Jest to \emph{ciągły gramian
obserwowalności na horyzoncie \(T\)}.
\end{definicja}

\begin{twierdzenie}[Ciągły gramian i odtworzenie stanu]
\label{thm:lin27-cont-gramian}
Dla każdego \(T>0\)
i \(x\in\mathbb C^n\)
zachodzi
\[
 x^*W_{\mathrm o,T}^{\mathrm c}x
 =\int_0^T\|Ce^{tA}x\|_2^2
  \,\mathrm dt.
\]
Macierz
\(W_{\mathrm o,T}^{\mathrm c}\)
jest hermitowska i nieujemnie
określona, a jej jądro jest
równe \(\mathcal N(C,A)\).
Para \((C,A)\) jest więc
obserwowalna dokładnie wtedy,
gdy \(W_{\mathrm o,T}^{\mathrm c}>0\)
dla pewnego, równoważnie
każdego \(T>0\).

Jeśli para jest obserwowalna
i znamy funkcję
\(y(t)=Ce^{tA}x_0\)
dla \(0\le t\le T\),
to stan początkowy
wyznacza wzór
\[
 x_0=(W_{\mathrm o,T}^{\mathrm c})^{-1}
       \int_0^T e^{tA^*}C^*y(t)
       \,\mathrm dt.
\]
\end{twierdzenie}
\begin{proof}
Przenosząc wektory poza całkę
i łącząc czynniki,
otrzymujemy
\[
 \begin{aligned}
 x^*W_{\mathrm o,T}^{\mathrm c}x
 &=\int_0^T
    (Ce^{tA}x)^*(Ce^{tA}x)
    \,\mathrm dt\\
 &=\int_0^T\|Ce^{tA}x\|_2^2
    \,\mathrm dt.
 \end{aligned}
\]
Hermitowskość i nieujemna
określoność wynikają stąd
bezpośrednio. Całka nieujemnej,
ciągłej funkcji jest zerowa
dokładnie wtedy, gdy
\(Ce^{tA}x=0\)
dla wszystkich \(t\in[0,T]\).
Różniczkując tę równość
w \(t=0\) dla
\(k=0,\ldots,n-1\),
otrzymujemy
\(CA^kx=0\).
Stąd \(x\in\mathcal N(C,A)\).
Odwrotnie, jeżeli
\(x\in\mathcal N(C,A)\),
to Cayley--Hamilton daje
\(CA^kx=0\) dla wszystkich
\(k\ge0\), a szereg
\(e^{tA}=\sum_{k\ge0}t^kA^k/k!\)
implikuje \(Ce^{tA}x=0\).
W ten sposób jądro gramianu
jest dokładnie \(\mathcal N(C,A)\)
i nie zależy od \(T>0\).

Dla obserwacji
\(y(t)=Ce^{tA}x_0\)
mamy
\[
 \int_0^T e^{tA^*}C^*y(t)
   \,\mathrm dt
 =\int_0^T e^{tA^*}C^*Ce^{tA}
   \,\mathrm dt\, x_0
 =W_{\mathrm o,T}^{\mathrm c}x_0.
\]
Ponieważ obserwowalność
zapewnia dodatnią określoność
gramianu, jest on odwracalny;
otrzymujemy wzór na \(x_0\).
\end{proof}

\begin{przyklad}[Jawne obliczenie obu gramianów]
\label{prz:lin27-two-gramians}
Niech \(n=2\), \(p=1\)
oraz
\[
 A=\begin{pmatrix}0&1\\0&0\end{pmatrix},
 \qquad C=\begin{pmatrix}1&0\end{pmatrix}.
\]
Wówczas \(CA=(0,1)\), zatem
\(\mathcal O_2=I_2\)
i
\(W_{\mathrm o,2}^{\mathrm d}=I_2\).

Ponieważ \(A^2=0\),
mamy \(e^{tA}=I+tA\) i
\(Ce^{tA}=(1,t)\).
Dlatego
\[
 W_{\mathrm o,T}^{\mathrm c}
 =\int_0^T
 \begin{pmatrix}1&t\\t&t^2\end{pmatrix}
 \,\mathrm dt
 =\begin{pmatrix}
    T&T^2/2\\
    T^2/2&T^3/3
   \end{pmatrix}.
\]
Jego wyznacznik to
\(T^4(1/3-1/4)=T^4/12>0\).
Stan można więc odtworzyć
na każdym dodatnim horyzoncie,
choć w bardzo krótkim czasie
odwracanie gramianu jest źle
uwarunkowane.
\end{przyklad}
